
Or d'après VI 4.8, @@M0@@ est la @@M1@@-algèbre
engendrée par les @@M2@@, soumis aux relations déduites de :
@@M3@@
Donc @@M4@@, où @@M5@@ est l'idéal engendré par les @@M6@@ pour @@M7@@, donc le @@M8@@-idéal engendré par @@M9@@. Par suite,
@@M10@@. On achève la démonstration par le lemme :

\textbf{Lemme 4.2.} \uline{Soient} @@M11@@ \uline{un schéma noethérien}, @@M12@@ \uline{un sous-groupe parabolique d'un} @@M13@@\uline{-groupe réductif}, @@M14@@ \uline{un torseur sous le radical unipotent} @@M15@@ \uline{de} @@M16@@; \uline{alors on a un isomorphisme} :
@@M17@@
En effet, il existe une suite de sous-schémas en groupes de @@M18@@ :
@@M19@@
fermés dans @@M20@@, finie car @@M21@@ est noethérien, et tels que, pour tout @@M22@@, @@M23@@ soit isomorphe à un @@M24@@-groupe vectoriel (SGA 3 XXVI 2.1). Par suite, @@M25@@ se factorise par une suite de torseurs sous des groupes vectoriels, d'où le résultat d'après 3.4.

\textbf{Corollaire 4.3.} \uline{Soient} @@M26@@ \uline{un schéma noethérien}, @@M27@@. \uline{Alors l'homomorphisme canonique} @@M28@@ \uline{est un isomorphisme}.

Résulte immédiatement de 4.1, appliqué à @@M29@@, puisqu'alors @@M30@@.

